%% =============================================================================
%% BEATRIX - Innosuisse Innovation Cheque Antragstext
%% Version 4.7 - Ernst-Fehr-Audit: BCM denkt/modelliert, BEATRIX rechnet/skaliert (11 Korrekturen)
%% Datum: 12. Februar 2026
%% Basis: V4.4 (27. Januar 2026) - alle Inhalte beibehalten
%% Änderung: Saubere 3-Ebenen-Terminologie: EBF (Evidenz) → BCM (Modell) → BEATRIX (Tool)
%% =============================================================================

\documentclass[11pt,a4paper]{article}

%% Packages
\usepackage[utf8]{inputenc}
\usepackage[T1]{fontenc}
\usepackage[ngerman]{babel}
\usepackage{geometry}
\usepackage{booktabs}
\usepackage{longtable}
\usepackage{array}
\usepackage{xcolor}
\usepackage{hyperref}
\usepackage{enumitem}
\usepackage{titlesec}
\usepackage{fancyhdr}
\usepackage{parskip}
\usepackage{microtype}
\usepackage{graphicx}

%% Page geometry
\geometry{
    left=2.5cm,
    right=2.5cm,
    top=2.5cm,
    bottom=2.5cm
}

%% Colors
\definecolor{fehrgray}{RGB}{80,80,80}
\definecolor{fehrblue}{RGB}{0,51,102}

%% Hyperref setup
\hypersetup{
    colorlinks=true,
    linkcolor=fehrblue,
    citecolor=fehrblue,
    urlcolor=fehrblue
}

%% Section formatting
\titleformat{\section}
    {\Large\bfseries\color{fehrblue}}
    {\thesection}{1em}{}
\titleformat{\subsection}
    {\large\bfseries\color{fehrgray}}
    {\thesubsection}{1em}{}
\titleformat{\subsubsection}
    {\normalsize\bfseries}
    {\thesubsubsection}{1em}{}

%% Header/Footer
\pagestyle{fancy}
\fancyhf{}
\fancyhead[L]{\small\textcolor{fehrgray}{BEATRIX -- Innovation Cheque Innosuisse}}
\fancyhead[R]{\small\textcolor{fehrgray}{SSBM}}
\fancyfoot[C]{\thepage}
\renewcommand{\headrulewidth}{0.4pt}

%% Table column types
\newcolumntype{L}[1]{>{\raggedright\arraybackslash}p{#1}}
\newcolumntype{C}[1]{>{\centering\arraybackslash}p{#1}}

\begin{document}

%% =============================================================================
%% TITELSEITE
%% =============================================================================

\begin{titlepage}
\centering
\vspace*{2cm}

{\Huge\bfseries\color{fehrblue} BEATRIX}\\[0.5cm]
{\Large Behavioral Economics AI Technology for\\Research and Implementation eXcellence}\\[2cm]

{\LARGE Antworten auf den\\Innosuisse Innovation Cheque Fragebogen}\\[1cm]

{\large Version 4.7 -- Scharfe Abgrenzung EBF / BCM / BEATRIX}\\[2cm]

\rule{0.8\textwidth}{0.5pt}\\[1.5cm]

\begin{tabular}{ll}
\textbf{Förderprogramm:} & Innosuisse Innovation Cheque \\[0.3cm]
\textbf{Fördersumme:} & CHF 15'000 \\[0.3cm]
\textbf{Fachgebiet:} & Social Sciences \& Business Management (SSBM) \\[0.3cm]
\textbf{Wirtschaftspartner:} & FehrAdvice Partners AG \\[0.3cm]
\textbf{Forschungspartner:} & Universität Zürich \\[0.3cm]
& Prof. Harald Gall \quad Institut für Informatik, UZH \\
& Prof. Johannes Luger \quad Institut für Betriebswirtschaft, UZH \\
& Prof. Ernst Fehr \quad Scientific Advisor \\
\end{tabular}

\vfill

\fbox{\parbox{0.9\textwidth}{\centering\large\bfseries
\textit{``Die Technologie ermöglicht die Innovation, aber die Transformation der\\Beratungspraxis ist das Produkt.''}
}}

\vfill

{\large Datum: 27. Januar 2026}

\end{titlepage}

%% =============================================================================
%% INHALTSVERZEICHNIS
%% =============================================================================

\tableofcontents
\newpage

%% =============================================================================
%% ABSTRACT
%% =============================================================================

\section*{Abstract}
\addcontentsline{toc}{section}{Abstract}

Unternehmen und öffentliche Verwaltungen verfügen über exzellente Datensysteme für harte Faktoren wie Finanzen, Logistik oder Produktion. Für die systematische Analyse von Verhaltenstreibern fehlen jedoch bis heute wissenschaftlich fundierte Werkzeuge. Die Folge ist bekannt: Rund 70\% aller Change-Projekte erreichen ihre Ziele nicht, weil menschliche Faktoren unterschätzt oder unzureichend berücksichtigt werden (Kotter, 1996; McKinsey, 2015).

FehrAdvice Partners verfügt über eine wissenschaftlich fundierte, über Jahre validierte methodische Kernkompetenz in der Analyse von Verhaltensentscheidungen. Diese Expertise wird durch das BCM (Behavioral Change Model) formalisiert und durch BEATRIX skalierbar gemacht. Das BCM überführt bisher personengebundene verhaltensökonomische Expertise in eine systematisch anwendbare Entscheidungslogik; BEATRIX setzt diese über unterschiedliche Kontexte hinweg konsistent ein.

BEATRIX ist ein KI-gestütztes Entscheidungsunterstützungssystem, das auf dem BCM (Behavioral Change Model) basiert. Das BCM formalisiert mehrere Jahrzehnte verhaltensökonomische Forschung; BEATRIX macht diese Formalisierung für die Managementpraxis systematisch zugänglich. Das BCM liefert drei Kernelemente: (1) systematische Kontextmodellierung, (2) die Erfassung von Wechselwirkungen, die in der Realität stattfinden (psychologische, finanzielle, kulturelle Effekte), sowie (3) eine transparente, nachvollziehbare Entscheidungslogik. BEATRIX operationalisiert diese Elemente und schlägt auf dieser Basis geeignete Massnahmen vor. Die Kombination verhaltensökonomischer Methodik mit KI ermöglicht es erstmals, mehrere relevante Faktoren simultan zu berücksichtigen und deren Interdependenzen systematisch in Entscheidungsprozesse zu integrieren.

Die Universität Zürich validiert BEATRIX mit zwei komplementären Kompetenzen: Prof. Harald Gall analysiert die technische Architektur im Hinblick auf Skalierbarkeit, während Prof. Johannes Luger die Nutzerakzeptanz untersucht und Einführungsstrategien entwickelt.

Der Innovation Cheque finanziert eine fokussierte Machbarkeitsstudie im Umfang von 125 Arbeitsstunden, bestehend aus einem Architektur-Assessment, einer Usability-Studie mit 10--15 Führungskräften sowie einem Validierungsbericht als Grundlage für ein Innosuisse-Hauptprojekt.

BEATRIX adressiert einen Markt von CHF 5--10 Milliarden im DACH-Raum (Organisations- und Verhaltensberatung). Pilotprojekte zeigen vielversprechende Resultate. Das Projekt stärkt die Schweizer Position als Standort für evidenzbasierte Management-Innovationen -- entwickelt in Zürich, für Schweizer KMU und weit darüber hinaus.

\vspace{0.5cm}
\noindent\textbf{Begriffliche Einordnung: EBF, BCM und BEATRIX.} Drei Begriffe durchziehen diesen Antrag:

\begin{itemize}[nosep]
    \item \textbf{EBF} (Evidence-Based Framework for Economic and Social Behavior): Die kuratierte wissenschaftliche Wissensbasis -- über 2'000 Studien, 153 Theorien und 852 dokumentierte Fälle, axiomatisch geordnet. Das EBF ist selbst \emph{kein} Prognosemodell, sondern die Evidenzgrundlage, aus der Modelle abgeleitet werden.
    \item \textbf{BCM} (Behavioral Change Model): Das formale Verhaltensmodell, abgeleitet aus dem EBF. Kernidee: Verhaltensparameter (z.\,B. Verlustaversion, Zeitpräferenz) sind keine Konstanten, sondern kontextabhängige Funktionen $\theta = f(\Psi)$. Das BCM formalisiert Kontext ($\Psi$-Vektor), Wechselwirkungen ($\gamma$-Matrix) und den Veränderungsprozess (Utility $\to$ Awareness $\to$ Willingness $\to$ Journey) -- und generiert damit testbare Vorhersagen.
    \item \textbf{BEATRIX} (Behavioral Economics AI Technology for Research and Implementation eXcellence): Das KI-gestützte Werkzeug, das das BCM operationalisiert. BEATRIX macht die BCM-Logik skalierbar, konsistent und für Nicht-Expert:innen zugänglich. Der Innovation Cheque validiert BEATRIX -- nicht EBF oder BCM.
\end{itemize}

\noindent\textit{Kurzformel: EBF liefert die Evidenz $\to$ BCM formalisiert das Modell $\to$ BEATRIX macht es anwendbar. Die Hierarchie ist strikt: ohne EBF kein BCM, ohne BCM kein BEATRIX.}

\newpage

%% =============================================================================
%% 1. INNOVATION POTENTIAL
%% =============================================================================

\section{Innovation Potential}

\subsection{Das Problem, das wir lösen}

Stellen Sie sich vor: Ein Unternehmen will die Mitarbeiterzufriedenheit verbessern. Der Berater empfiehlt zehn Massnahmen -- aber welche davon wirken zusammen? Welche blockieren sich gegenseitig? Und funktioniert das, was in Skandinavien klappt, auch in der Schweiz?

\textbf{Heute basiert diese Einschätzung auf Erfahrung und Intuition. BEATRIX macht sie systematisch und nachvollziehbar.}

BEATRIX geht über Standardberatung hinaus, indem es auf Basis des BCM Verhaltensentscheidungen nicht erfahrungsbasiert beurteilt, sondern Kontext ($\Psi$), Wechselwirkungen ($\gamma$) und Entscheidungslogik systematisch modelliert und nachvollziehbar macht.

\subsection{Die drei Säulen der Innovation}

\subsubsection{Säule 1: Kontext ernst nehmen}

Menschen verhalten sich nicht isoliert vom Kontext (Fehr \& Gächter, 2002a). Ein finanzieller Anreiz wirkt anders bei einem Berufseinsteiger als bei einer Führungskraft kurz vor der Pensionierung (Gneezy et al., 2011). \textbf{BEATRIX nutzt das BCM, um acht Kontext-Dimensionen systematisch zu berücksichtigen} -- formalisiert als $\Psi$-Vektor (vgl. Kapitel 9):

\begin{table}[h]
\centering
\begin{tabular}{@{}L{4cm}L{9cm}@{}}
\toprule
\textbf{Dimension} & \textbf{Frage} \\
\midrule
Zeitlicher Kontext & Ist die Entscheidung dringend oder langfristig? \\
Sozialer Kontext & Wer beobachtet? Wer urteilt? \\
Institutioneller Kontext & Welche Regeln gelten? Welche Kultur herrscht? \\
Informationeller Kontext & Was weiss die Person? Was fehlt? \\
Emotionaler Kontext & Ist die Person gestresst? Entspannt? \\
Ressourcen-Kontext & Hat sie Zeit, Geld, Energie für die Entscheidung? \\
Kultureller Kontext & Schweiz ist nicht Schweden ist nicht Singapur \\
Physischer Kontext & Wo findet die Entscheidung statt? \\
\bottomrule
\end{tabular}
\end{table}

\textbf{Warum ist das neu?} Bisherige Ansätze behandeln Kontext als Störfaktor. Das BCM macht ihn zum zentralen Ordnungsprinzip -- BEATRIX setzt dies operativ um (Thaler \& Sunstein, 2008).

\subsubsection{Säule 2: Wechselwirkungen in der Realität erfassen -- und Massnahmen vorschlagen}

BEATRIX analysiert nicht Massnahmen isoliert, sondern die Realität dahinter: \textbf{Das BCM erfasst Wechselwirkungen, die in der Realität selbst stattfinden} (psychologisch, finanziell, kulturell) -- und BEATRIX schlägt auf dieser Basis geeignete Massnahmen vor. Diese realen Wechselwirkungen sind extrem schwierig zu erfassen und umfassen drei Ebenen:

\begin{table}[h]
\centering
\begin{tabular}{@{}L{4cm}L{9cm}@{}}
\toprule
\textbf{Wechselwirkungs-Ebene} & \textbf{Was das BCM modelliert} \\
\midrule
Psychologische Effekte & Wie wirken Anreize auf Motivation, Fairness-Wahrnehmung, intrinsische vs. extrinsische Motivation? (Fehr \& Schmidt, 1999; Deci et al., 1999) \\
Finanzielle Effekte & Wie interagieren monetäre Anreize mit nicht-monetären Faktoren? Wann verdrängen Boni intrinsische Motivation? (Frey \& Oberholzer-Gee, 1997; Gneezy \& Rustichini, 2000) \\
Kultureller Kontext & Wie variieren diese Effekte zwischen Kulturen? Was in Skandinavien funktioniert, kann in der Schweiz anders wirken. \\
\bottomrule
\end{tabular}
\end{table}

Auf Basis dieser Realitäts-Analyse identifiziert BEATRIX, welche Massnahmen sich verstärken oder blockieren:

\begin{table}[h]
\centering
\begin{tabular}{@{}L{3cm}L{4cm}L{6cm}@{}}
\toprule
\textbf{Kombination} & \textbf{Effekt} & \textbf{Beispiel} \\
\midrule
Verstärkend & Wirken zusammen stärker & Transparenz + Feedbackkultur \\
Neutralisierend & Heben sich auf & Individualboni + Teamziele \\
Blockierend & Eine verhindert die andere & Kontrolle + Vertrauenskultur \\
\bottomrule
\end{tabular}
\end{table}

\textbf{Warum ist das neu?} Traditionelle Beratung betrachtet Massnahmen isoliert. Das BCM modelliert die komplexen Wechselwirkungen, die in der Realität zwischen psychologischen, finanziellen und kulturellen Faktoren stattfinden -- BEATRIX \textbf{schlägt auf dieser Basis passende Massnahmen vor}.

\subsubsection{Säule 3: Transparente Entscheidungslogik}

BEATRIX liefert keine undurchsichtigen Empfehlungen. Jede Analyse ist nachvollziehbar und basiert auf dokumentierter Forschung.

\begin{table}[h]
\centering
\begin{tabular}{@{}L{6cm}L{7cm}@{}}
\toprule
\textbf{Traditionell} & \textbf{Mit BEATRIX} \\
\midrule
``Der Berater sagt, wir sollen X tun.'' & ``Die Analyse zeigt: In unserem Kontext wirkt X, weil Y.'' \\
Vertrauen in Person & Vertrauen in Methode \\
Schwer replizierbar & Nachvollziehbar und lernend \\
\bottomrule
\end{tabular}
\end{table}

\subsection{Was BEATRIX nicht ist}

Wir sind ehrlich über die Grenzen:

\begin{itemize}
    \item \textbf{Kein Wahrsager:} Das BCM modelliert Verhaltensparameter auf Populationsebene -- BEATRIX sagt nicht vorher, was eine einzelne Person tun wird
    \item \textbf{Kein Ersatz für Führung:} Das System unterstützt Entscheidungen, es trifft sie nicht
    \item \textbf{Kein Allheilmittel:} Manche Probleme sind nicht durch bessere Analyse lösbar
    \item \textbf{Keine ``künstliche Intelligenz'' im Hype-Sinn:} Systematische Analyse, kein magisches Orakel
\end{itemize}

\subsection{Was validiert der Forschungspartner?}

\begin{table}[h]
\centering
\begin{tabular}{@{}L{3cm}L{5cm}L{5cm}@{}}
\toprule
\textbf{Wer} & \textbf{Fragt} & \textbf{Methode} \\
\midrule
Prof. Gall & Ist das System robust und skalierbar? & Technische Analyse \\
Prof. Luger & Vertrauen Manager den Empfehlungen? & Tests mit Führungskräften \\
\bottomrule
\end{tabular}
\end{table}

\vspace{1cm}
\fbox{\parbox{0.95\textwidth}{
\textbf{Zusammenfassung: Innovation Potential}

BEATRIX ist eine Management-Innovation, keine Software-Innovation. Das System weiterentwickelt die Beratungspraxis systematisch, indem es:
\begin{enumerate}
    \item Kontext als zentrales Ordnungsprinzip verwendet
    \item Wechselwirkungen erfasst, die in der Realität stattfinden (psychologische, finanzielle, kulturelle Effekte) -- und darauf basierend Massnahmen vorschlägt
    \item Transparente Entscheidungslogik bietet statt undurchsichtiger Empfehlungen
\end{enumerate}
Die UZH prüft: Ist das System technisch robust (Gall)? Vertrauen Manager den Empfehlungen (Luger)?
}}

\newpage

%% =============================================================================
%% 2. MARKET RELEVANCE
%% =============================================================================

\section{Market Relevance}

\subsection{Die Soft-Factor-Lücke}

Moderne Unternehmen haben für fast alles Daten:

\begin{table}[h]
\centering
\begin{tabular}{@{}L{4cm}L{4cm}L{5cm}@{}}
\toprule
\textbf{Bereich} & \textbf{Datenqualität} & \textbf{Werkzeuge} \\
\midrule
Finanzen & Exzellent & SAP, Oracle \\
Logistik & Sehr gut & Supply-Chain-Systeme \\
Marketing & Gut & CRM, Analytics \\
Mitarbeiterverhalten & Lückenhaft & Umfragen, Intuition \\
Kundenentscheidungen & Lückenhaft & Fokusgruppen, Bauchgefühl \\
\bottomrule
\end{tabular}
\end{table}

\textbf{Das Problem:} Für die ``harten'' Faktoren gibt es hervorragende Systeme. Für die ``weichen'' Faktoren -- menschliches Verhalten, Motivation, Entscheidungen -- fehlen systematische Werkzeuge.

\subsection{Die Konsequenz}

Studien zeigen: \textbf{Etwa 70\% aller Change-Projekte erreichen ihre Ziele nicht} (Kotter, 1996; McKinsey, 2015). Der häufigste Grund? Menschliche Faktoren wurden unterschätzt.

\begin{table}[h]
\centering
\begin{tabular}{@{}L{4cm}L{4cm}L{5cm}@{}}
\toprule
\textbf{Typisches Szenario} & \textbf{Was passiert} & \textbf{Mit systematischer Analyse} \\
\midrule
Neue Software einführen & Mitarbeiter nutzen sie nicht & Widerstände früh erkennen \\
Kultur verändern & Alte Gewohnheiten bleiben & Hebel identifizieren \\
Anreizsystem ändern & Unerwartete Nebenwirkungen & Wechselwirkungen prüfen \\
\bottomrule
\end{tabular}
\end{table}

\subsection{Für wen ist BEATRIX gedacht?}

BEATRIX ist ein Tool zur Entscheidungsunterstützung für alle Entscheidungsträger*innen in Organisationen und Unternehmen, die grundlegende Entscheidungen treffen müssen, die von vielen gleichzeitig miteinander verflochtenen Faktoren geprägt sind (also von komplexen Entscheidungen). Das umfasst beispielsweise:

\begin{table}[h]
\centering
\begin{tabular}{@{}L{3.5cm}L{4.5cm}L{5cm}@{}}
\toprule
\textbf{Wer?} & \textbf{Typische Frage} & \textbf{Wie BEATRIX hilft} \\
\midrule
Beratungsunternehmen & ``Wie können wir unsere Empfehlungen besser begründen?'' & Systematische Analyse statt Intuition \\
Personalabteilungen & ``Warum funktioniert unsere neue Strategie nicht?'' & Kontext-Analyse der Belegschaft \\
Öffentliche Verwaltung & ``Wie erreichen wir, dass Bürger das neue Angebot nutzen?'' & Evidenzbasierte Gestaltung \\
Geschäftsführung & ``Welche Massnahmen wirken wirklich?'' & Priorisierung nach Wirksamkeit \\
\bottomrule
\end{tabular}
\end{table}

\subsection{Ein konkretes Beispiel}

\textbf{Situation:} Ein Energieversorger will Kunden zum Wechsel auf einen günstigeren Tarif bewegen.

\textbf{Traditioneller Ansatz:} Brief schicken, hoffen, dass Kunden reagieren.

\textbf{Mit systematischer Analyse:}
\begin{itemize}
    \item Kontext verstehen: Warum wechseln Kunden nicht, obwohl sie sparen würden? (Kahneman \& Tversky, 1979)
    \item Hindernisse identifizieren: Komplexität? Trägheit? Misstrauen?
    \item Massnahmen anpassen: Einfacherer Prozess? Persönliche Ansprache? (Thaler \& Sunstein, 2008)
\end{itemize}

\textbf{Ergebnis in Pilotprojekten:} Deutlich höhere Wechselquoten -- weil die Hindernisse verstanden wurden.

\subsection{Marktgrösse}

\begin{table}[h]
\centering
\begin{tabular}{@{}L{4cm}L{6cm}L{3cm}@{}}
\toprule
\textbf{Ebene} & \textbf{Beschreibung} & \textbf{Schätzung} \\
\midrule
Gesamtmarkt (TAM) & Beratungsmarkt weltweit & USD 300+ Mrd. \\
Relevanter Markt (SAM) & Organisations-/Verhaltensberatung (DACH) & CHF 5--10 Mrd. \\
Erreichbarer Markt (SOM) & Was wir in 5 Jahren realistisch erreichen können & CHF 50--100 Mio. \\
\bottomrule
\end{tabular}
\end{table}

\textbf{Hinweis:} Diese Schätzungen sind konservativ und hängen vom Projekterfolg ab.

\vspace{1cm}
\fbox{\parbox{0.95\textwidth}{
\textbf{Zusammenfassung: Market Relevance}

BEATRIX adressiert eine echte Lücke: Während für harte Faktoren exzellente Werkzeuge existieren, fehlt für Verhaltensfragen ein systematischer Ansatz.
\begin{itemize}
    \item \textbf{Klares Problem:} Hohe Scheiterquote bei Veränderungsprojekten
    \item \textbf{Definierte Zielgruppen:} Beratungen, Personalabteilungen, öffentliche Verwaltung
    \item \textbf{Validierte Wirksamkeit:} Positive Ergebnisse in Pilotprojekten
    \item \textbf{Realistisches Marktpotenzial:} Konservativ geschätzt
\end{itemize}
}}

\newpage

%% =============================================================================
%% 3. FEASIBILITY & RISK REDUCTION
%% =============================================================================

\section{Feasibility \& Risk Reduction}

\subsection{Was wurde bereits entwickelt?}

\begin{table}[h]
\centering
\begin{tabular}{@{}L{3cm}L{3cm}L{7cm}@{}}
\toprule
\textbf{Phase} & \textbf{Zeitraum} & \textbf{Was wurde erreicht?} \\
\midrule
Konzeption & 2018--2020 & Theoretisches Fundament erarbeitet \\
Prototyp & 2020--2022 & Erste funktionierende Version \\
Pilotphase & 2022--2025 & Einsatz bei über 20 Kundenprojekten \\
Heute & 2026 & Funktionierender Kern, bereit für Skalierung \\
\bottomrule
\end{tabular}
\end{table}

\textbf{Aktueller Stand:} Das System funktioniert und wird eingesetzt. Was fehlt, ist die externe Validierung der Skalierbarkeit und der Nutzerakzeptanz.

\subsection{Was funktioniert bereits?}

\begin{itemize}
    \item[$\checkmark$] Systematische Erfassung der acht Kontext-Dimensionen
    \item[$\checkmark$] Erfassung von Wechselwirkungen, die in der Realität stattfinden, und darauf basierende Massnahmenvorschläge
    \item[$\checkmark$] Benutzeroberfläche für Berater
    \item[$\checkmark$] Bibliothek mit über 50 dokumentierten Anwendungsfällen
    \item[$\checkmark$] Integration von über 1'500 wissenschaftlichen Studien
\end{itemize}

\subsection{Die zwei offenen Fragen}

\begin{table}[h]
\centering
\begin{tabular}{@{}L{5cm}L{5cm}L{3cm}@{}}
\toprule
\textbf{Frage} & \textbf{Warum kritisch?} & \textbf{Wer prüft?} \\
\midrule
Ist das System skalierbar? & Ohne Skalierbarkeit kein Geschäftsmodell & Prof. Gall \\
Vertrauen Manager den Empfehlungen? & Ohne Akzeptanz keine Nutzung & Prof. Luger \\
\bottomrule
\end{tabular}
\end{table}

\subsection{Was kann mit CHF 15'000 erreicht werden?}

Der Innovation Cheque finanziert 125 Arbeitsstunden (à CHF 120):

\begin{table}[h]
\centering
\begin{tabular}{@{}L{4cm}L{3cm}L{6cm}@{}}
\toprule
\textbf{Arbeitspaket} & \textbf{Umfang} & \textbf{Ergebnis} \\
\midrule
Architektur-Analyse & 40 Stunden & Technischer Prüfbericht \\
Nutzer-Tests & 50 Stunden & Akzeptanz-Studie \\
Risikoanalyse & 20 Stunden & Identifikation kritischer Punkte \\
Abschlussbericht & 15 Stunden & Empfehlungen für Hauptprojekt \\
\midrule
\textbf{Total} & \textbf{125 Stunden} & \\
\bottomrule
\end{tabular}
\end{table}

\subsection{Entwicklungspfad}

\begin{center}
\begin{tabular}{c}
\textbf{Heute (2026)} $\rightarrow$ Wir sind hier \\
$\downarrow$ \\
\textbf{Innovation Cheque} -- Machbarkeit prüfen (CHF 15k) \\
$\downarrow$ \\
\textbf{Hauptprojekt} -- Skalierbare Lösung (CHF 500k--2M) \\
$\downarrow$ \\
\textbf{Marktreife (2028--29)} -- Breiter Einsatz möglich \\
\end{tabular}
\end{center}

\subsection{Entscheidungskriterien nach Innovation Cheque}

Am Ende der Machbarkeitsstudie wird eine fundierte Entscheidung getroffen:

\begin{table}[h]
\centering
\begin{tabular}{@{}L{2cm}L{7cm}L{4cm}@{}}
\toprule
\textbf{Szenario} & \textbf{Kriterium} & \textbf{Entscheidung} \\
\midrule
GO & Architektur skalierbar, Akzeptanz bei $\geq$70\% der Testpersonen & Hauptprojekt beantragen \\
PILOT & Teilweise Anpassungen nötig, Akzeptanz 50--70\% & Fokussiertes Folgeprojekt \\
NO-GO & Fundamentale Architektur-Probleme oder Akzeptanz <50\% & Projekt nicht weiterverfolgen \\
\bottomrule
\end{tabular}
\end{table}

Diese transparenten Kriterien stellen sicher, dass die Ergebnisse des Innovation Cheque zu einer klaren, evidenzbasierten Entscheidung führen.

\vspace{1cm}
\fbox{\parbox{0.95\textwidth}{
\textbf{Zusammenfassung: Feasibility \& Risk Reduction}

BEATRIX ist keine Idee auf dem Papier -- das System funktioniert und wird eingesetzt.

Der Innovation Cheque klärt zwei kritische Fragen:
\begin{itemize}
    \item Ist das System technisch bereit für breiteren Einsatz? (Gall)
    \item Vertrauen Führungskräfte den Empfehlungen? (Luger)
\end{itemize}

Mit CHF 15'000 werden 125 Arbeitsstunden (à CHF 120) finanziert, die eine fundierte Go/Pilot/No-Go-Entscheidung für das Hauptprojekt ermöglichen.
}}

\newpage

%% =============================================================================
%% 4. COLLABORATION WITH RESEARCH PARTNER
%% =============================================================================

\section{Collaboration with Research Partner}

\subsection{Was FehrAdvice mitbringt -- und was fehlt}

\begin{table}[h]
\centering
\begin{tabular}{@{}L{6cm}L{7cm}@{}}
\toprule
\textbf{Stärke von FehrAdvice} & \textbf{Was fehlt} \\
\midrule
Verhaltensökonomie (15+ Jahre) & Technische Architektur für Skalierung \\
Praxiswissen (20+ Projekte) & Wissenschaftliche Akzeptanzforschung \\
Kundenzugang & Forschungsbasiertes Change Management \\
Ernst Fehr als Gründer & Unabhängige externe Prüfung \\
\bottomrule
\end{tabular}
\end{table}

\subsection{Warum die Universität Zürich?}

\begin{table}[h]
\centering
\begin{tabular}{@{}L{4cm}L{4cm}L{5cm}@{}}
\toprule
\textbf{Kriterium} & \textbf{Anforderung} & \textbf{UZH} \\
\midrule
Wissenschaftliche Qualität & Höchste Standards & $\checkmark$ International anerkannt \\
Thematische Passung & Software + Management & $\checkmark$ Gall + Luger \\
Bestehende Beziehung & Vertrauensbasis & $\checkmark$ Fehr ist UZH-Professor \\
Nähe & Enge Zusammenarbeit möglich & $\checkmark$ Alle in Zürich \\
\bottomrule
\end{tabular}
\end{table}

\subsection{Die zwei Professoren}

\textbf{Prof. Harald Gall (Institut für Informatik)}

Renommierter Experte für Software-Entwicklung mit über 300 Publikationen. Seine Frage:

\textit{``Wie muss ein System wie BEATRIX gebaut sein, damit es zuverlässig wächst und sich weiterentwickelt?''}

\textbf{Prof. Johannes Luger (Institut für Betriebswirtschaft)}

Forscht zur Zusammenarbeit zwischen Menschen und intelligenten Systemen -- zentral für die Willingness-Komponente des BCM (Kapitel 12) (Krakowski et al., 2025). Seine Frage:

\textit{``Unter welchen Bedingungen vertrauen Führungskräfte den Empfehlungen eines Analyse-Systems?''}

\subsection{Konkrete Beiträge}

\begin{table}[h]
\centering
\begin{tabular}{@{}L{3cm}L{5cm}L{5cm}@{}}
\toprule
\textbf{Wer} & \textbf{Beitrag} & \textbf{Ergebnis} \\
\midrule
Prof. Gall & Architektur-Analyse & Prüfbericht, Skalierungs-Empfehlungen \\
Prof. Luger & Nutzer-Tests mit 10--15 Führungskräften & Akzeptanz-Studie, Einführungsstrategien \\
Beide & Risikoanalyse & Validierungsbericht für Hauptprojekt \\
\bottomrule
\end{tabular}
\end{table}

\vspace{1cm}
\fbox{\parbox{0.95\textwidth}{
\textbf{Zusammenfassung: Research Partner Collaboration}

Die UZH ist der ideale Partner für dieses Projekt:
\begin{itemize}
    \item Ernst Fehr verbindet FehrAdvice und UZH bereits
    \item Prof. Gall sichert die technische Qualität
    \item Prof. Luger erforscht die menschlichen Faktoren
    \item Beide bringen unabhängige wissenschaftliche Prüfung
\end{itemize}
Die Partnerschaft schliesst kritische Kompetenzlücken und legt den Grundstein für langfristige Zusammenarbeit.
}}

\newpage

%% =============================================================================
%% 5. IMPACT FOR SWITZERLAND
%% =============================================================================

\section{Impact for Switzerland}

\subsection{Eine Schweizer Lösung}

\begin{table}[h]
\centering
\begin{tabular}{@{}L{5cm}L{8cm}@{}}
\toprule
\textbf{Aspekt} & \textbf{Bedeutung} \\
\midrule
Geistiges Eigentum & bleibt in der Schweiz \\
Daten & unter Schweizer Recht \\
Lokale Wertschöpfung & Arbeitsplätze, Steuern, Know-how in der Schweiz \\
``Made in Switzerland'' & Qualitätsmerkmal auch für Beratungsdienstleistungen \\
\bottomrule
\end{tabular}
\end{table}

\subsection{Arbeitsplätze und Wertschöpfung}

\begin{table}[h]
\centering
\begin{tabular}{@{}L{5cm}L{8cm}@{}}
\toprule
\textbf{Zeitraum} & \textbf{Erwartung} \\
\midrule
2026 (Machbarkeitsphase) & 2--3 direkte Arbeitsplätze \\
2028 (Markteinführung) & 10--15 direkte Arbeitsplätze \\
2030 (Skalierung) & 30--50 direkte Arbeitsplätze \\
\bottomrule
\end{tabular}
\end{table}

\textbf{Hinweis:} Diese Schätzungen sind konservativ und hängen vom Projekterfolg ab.

\subsection{Die differenzierte Schweizer Position}

\begin{table}[h]
\centering
\begin{tabular}{@{}L{3cm}L{10cm}@{}}
\toprule
\textbf{Faktor} & \textbf{Schweiz-Vorteil} \\
\midrule
Ernst Fehr & Hochzitierter Verhaltensökonom (h-Index > 150) arbeitet hier \\
UZH-Forschung & Hervorragend positioniert in relevanten Feldern \\
FehrAdvice & Einziges Unternehmen dieser Art mit direkter Verbindung zur Spitzenforschung \\
Neutralität & Schweiz als vertrauenswürdiger Standort für sensible Analysen \\
Stabilität & Verlässliches Umfeld für langfristige Entwicklung \\
\bottomrule
\end{tabular}
\end{table}

\subsection{Gesellschaftlicher Nutzen}

Wenn Organisationen systematischer mit Verhaltensfragen umgehen, profitiert die Gesellschaft (Thaler \& Sunstein, 2008; Sunstein, 2014):

\begin{table}[h]
\centering
\begin{tabular}{@{}L{3cm}L{5cm}L{5cm}@{}}
\toprule
\textbf{Bereich} & \textbf{Heute} & \textbf{Mit systematischem Ansatz} \\
\midrule
Gesundheit & Präventionsprogramme werden wenig genutzt & Bessere Gestaltung erhöht Teilnahme \\
Umwelt & Appelle verhallen & Verhaltensbasierte Anreize wirken \\
Finanzen & Viele sparen zu wenig & Bessere Defaults helfen \\
Verwaltung & Angebote werden nicht genutzt & Bürgerfreundliche Gestaltung \\
\bottomrule
\end{tabular}
\end{table}

\vspace{1cm}
\fbox{\parbox{0.95\textwidth}{
\textbf{Zusammenfassung: Impact for Switzerland}

BEATRIX könnte der Schweizer Wirtschaft auf mehreren Ebenen nützen:
\begin{itemize}
    \item \textbf{Arbeitsplätze:} Konservativ geschätzt 30--50 direkt bis 2030
    \item \textbf{Know-how:} Schweizer Expertise in zukunftsrelevantem Feld
    \item \textbf{Modellcharakter:} Vorlage für KMU-Hochschul-Zusammenarbeit
    \item \textbf{Gesellschaft:} Bessere Entscheidungen in Gesundheit, Umwelt, Verwaltung
\end{itemize}
Die Schweiz hat eine differenzierte Ausgangslage durch die Kombination von Ernst Fehr, UZH und FehrAdvice. Dieses Potenzial verantwortungsvoll zu nutzen, ist das Ziel dieses Projekts.
}}

\newpage

%% =============================================================================
%% GESAMTZUSAMMENFASSUNG
%% =============================================================================

\section*{Gesamtzusammenfassung}
\addcontentsline{toc}{section}{Gesamtzusammenfassung}

\begin{table}[h]
\centering
\begin{tabular}{@{}L{3cm}L{10cm}@{}}
\toprule
\textbf{Frage} & \textbf{Kernaussage} \\
\midrule
a) Innovation & BEATRIX nutzt das BCM, um Wechselwirkungen zu modellieren, die in der Realität stattfinden (psychologisch, finanziell, kulturell) -- und schlägt darauf basierend Massnahmen vor. Systematisch und nachvollziehbar. \\
b) Markt & Die ``Soft-Factor-Lücke'' ist real: Für Finanzen gibt es SAP, für Verhalten fehlen systematische Werkzeuge \\
c) Machbarkeit & Das System funktioniert bereits -- der Innovation Cheque prüft Skalierbarkeit und Nutzerakzeptanz \\
d) Partnerschaft & UZH bringt genau die Kompetenzen, die FehrAdvice fehlen: technische Architektur und Akzeptanzforschung \\
e) Schweiz & Differenzierte Ausgangslage durch Fehr + UZH + FehrAdvice -- verantwortungsvoll nutzen \\
\bottomrule
\end{tabular}
\end{table}

\vspace{1cm}

\begin{center}
\fbox{\parbox{0.9\textwidth}{\centering
\textbf{Der strategische Kernsatz}

\vspace{0.5cm}

\textit{``Dieses Projekt entwickelt keine neue Software und keine neue Methodik -- das BCM existiert bereits. Es validiert die KI-Implementierung (BEATRIX) einer wissenschaftlich fundierten Methodik für bessere Entscheidungen in Organisationen. BEATRIX macht diese Methodik skalierbar -- die Transformation der Beratungspraxis ist das eigentliche Produkt.''}
}}
\end{center}

\newpage

%% =============================================================================
%% GLOSSAR
%% =============================================================================

\section*{Glossar}
\addcontentsline{toc}{section}{Glossar}

\begin{longtable}{@{}L{4cm}L{10cm}@{}}
\toprule
\textbf{Begriff} & \textbf{Definition} \\
\midrule
\endfirsthead
\toprule
\textbf{Begriff} & \textbf{Definition} \\
\midrule
\endhead

Awareness & Das Wissen einer Person darüber, dass ein Problem existiert und dass Handlungsoptionen verfügbar sind. Erste Stufe vor jeder Verhaltensänderung. Ohne Awareness keine Handlung -- aber Awareness allein reicht nicht aus (siehe Willingness). \\
\midrule
BEATRIX & Behavioral Economics AI Technology for Research and Implementation eXcellence. Ein KI-gestütztes Entscheidungsunterstützungssystem, das auf dem BCM (Behavioral Change Model) basiert. BEATRIX operationalisiert das BCM und macht damit 50 Jahre verhaltensökonomische Forschung (Kahneman, Thaler, Fehr) skalierbar für die Managementpraxis. Erfasst Wechselwirkungen, die in der Realität stattfinden (psychologische, finanzielle, kulturelle Effekte) -- und schlägt auf dieser Basis geeignete Massnahmen vor. \\
\midrule
Behavioral Change Journey (BCJ) & Der Weg, den eine Person von der ersten Wahrnehmung eines Problems bis zur dauerhaften Verhaltensänderung zurücklegt. Verläuft in Phasen: Aufmerksamkeit $\leftrightarrow$ Verstehen $\leftrightarrow$ Absicht $\leftrightarrow$ Handlung $\leftrightarrow$ Gewohnheit. Jede Phase erfordert andere Interventionen. \\
\midrule
BCM & Behavioral Change Model. Das formale Verhaltensmodell, das aus dem EBF (Evidence-Based Framework) abgeleitet ist. Formalisiert, wie Kontext ($\Psi$), Wechselwirkungen ($\gamma$) und der Veränderungsprozess (Utility $\rightarrow$ Awareness $\rightarrow$ Willingness $\rightarrow$ Journey) zusammenwirken. Das BCM ist die theoretische Grundlage von BEATRIX. \\
\midrule
Bias & Systematische Abweichung vom ``rationalen'' Entscheiden. Nicht zufällige Fehler, sondern vorhersagbare Muster. Beispiel: Wir gewichten Verluste etwa doppelt so stark wie gleichgrosse Gewinne (Verlustaversion). \\
\midrule
Change-Projekt & Geplante organisatorische Veränderung wie Prozesseinführung, Kulturwandel oder Restrukturierung. Etwa 70\% scheitern, weil menschliche Faktoren unterschätzt werden (Kotter, 1996). BEATRIX hilft, diese Faktoren systematisch zu analysieren. \\
\midrule
Choice Architecture & Die Gestaltung der Umgebung, in der Entscheidungen getroffen werden. Beeinflusst Verhalten, ohne Optionen einzuschränken. Beispiel: Obst auf Augenhöhe in der Kantine platzieren statt Süssigkeiten. \\
\midrule
Crowding-Out & Verdrängung intrinsischer Motivation durch externe Anreize. Klassisches Beispiel: Bezahlung für Blutspenden senkt die Spendenbereitschaft, weil die soziale Norm (``Helfen'') durch eine Marktnorm (``Verkaufen'') ersetzt wird (Frey \& Oberholzer-Gee, 1997). \\
\midrule
DACH-Raum & Wirtschaftsraum Deutschland, Österreich, Schweiz. Etwa 100 Millionen Einwohner, ähnliche Sprache und Geschäftskultur, aber unterschiedliche institutionelle Kontexte. \\
\midrule
Default & Voreingestellte Option, die gilt, wenn keine aktive Entscheidung getroffen wird. Wirkungsvoll, weil Menschen bei Unsicherheit oder Trägheit beim Status quo bleiben. Beispiel: Automatische Pensionskassen-Anmeldung erhöht Sparquoten um 30--50\% (Thaler \& Sunstein, 2008). \\
\midrule
Dual-System-Theorie & Modell von Kahneman: Menschen denken in zwei Modi. System 1: schnell, automatisch, intuitiv, fehleranfällig. System 2: langsam, bewusst, analytisch, anstrengend. Die meisten Alltagsentscheidungen trifft System 1 (Kahneman, 2011). \\
\midrule
EBF-Framework & Evidence-Based Framework for Economic and Social Behavior. Die eigenständige wissenschaftliche Wissensbasis: über 2'000 Studien, 153 Theorien, 852 dokumentierte Fälle, systematisiert in 21 Kapiteln, 56 Appendices und 414 formalen Axiomen. Das EBF existiert unabhängig von BEATRIX und integriert die Forschung von Kahneman, Thaler, Fehr und anderen in ein einheitliches System mit dem 10C-CORE. Aus dem EBF ist das BCM (Behavioral Change Model) abgeleitet, das BEATRIX operationalisiert. \\
\midrule
Evidenzbasiert & Auf wissenschaftlichen Belegen beruhend, nicht auf Intuition oder Erfahrung allein. Bedeutet: Aussagen sind durch Studien, Experimente oder systematische Beobachtungen gestützt. BEATRIX integriert über 1'500 wissenschaftliche Studien. \\
\midrule
Extrinsische Motivation & Antrieb durch äussere Faktoren wie Geld, Lob, Strafen oder soziale Anerkennung. Gegenteil: intrinsische Motivation. Beide können sich gegenseitig verstärken oder verdrängen (Deci et al., 1999). \\
\midrule
Fehr, Ernst & Schweizer Ökonom, Professor an der Universität Zürich, Mitgründer von FehrAdvice. Einer der meistzitierten Wirtschaftswissenschaftler weltweit. Forschungsschwerpunkte: Fairness, Kooperation, Bestrafungsverhalten (Fehr \& Schmidt, 1999; Fehr \& Gächter, 2000). Scientific Advisor von BEATRIX. \\
\midrule
FEPSDE & Die sechs Dimensionen von Wohlbefinden im BCM: Financial, Experiential, Procedural, Social, Distributive, Episodic. Jede Intervention wirkt auf eine oder mehrere dieser Dimensionen. \\
\midrule
Framing & Die Art, wie Information präsentiert wird, beeinflusst Entscheidungen. ``90\% Überlebensrate'' klingt besser als ``10\% Sterberate'' -- obwohl es dasselbe bedeutet. Framing ist kein Betrug, sondern unvermeidlich: Jede Aussage hat einen Rahmen. \\
\midrule
GO/PILOT/NO-GO & Entscheidungsmatrix nach Machbarkeitsstudien. GO: Kriterien erfüllt, Hauptprojekt starten. PILOT: Teilweise erfüllt, fokussiertes Folgeprojekt. NO-GO: Fundamentale Probleme, Projekt beenden. Verhindert Sunk-Cost-Fallen durch vordefinierte Abbruchkriterien. \\
\midrule
Heuristik & Mentale Abkürzung, die schnelle Entscheidungen ermöglicht. Oft nützlich, manchmal irreführend. Beispiel: ``Was mir leicht einfällt, ist wichtig'' (Verfügbarkeitsheuristik). Heuristiken sind keine Dummheit, sondern evolutionär sinnvolle Strategien mit Nebenwirkungen. \\
\midrule
Innovation Cheque & Innosuisse-Förderinstrument für Machbarkeitsstudien. CHF 15'000 für 125 Arbeitsstunden Forschungsleistung. Ziel: Validierung einer Idee vor grösseren Investitionen. Niedrige Hürde, um KMU den Zugang zu Hochschulforschung zu erleichtern. \\
\midrule
Innosuisse & Schweizerische Agentur für Innovationsförderung. Finanziert Zusammenarbeit zwischen Unternehmen und Forschungsinstitutionen. Budget: etwa CHF 500 Millionen jährlich. \\
\midrule
Intervention & Gezielte Massnahme zur Beeinflussung von Verhalten. Im BCM in acht Grundtypen klassifiziert: Awareness, Feedback, Kontextgestaltung, Timing, Identität, Soziale Normen, Finanzielle Anreize, Selbstverpflichtung. \\
\midrule
Intrinsische Motivation & Antrieb durch die Tätigkeit selbst -- Freude, Interesse, Sinnerleben. Stabiler und nachhaltiger als extrinsische Motivation, aber schwerer von aussen zu beeinflussen. Kann durch falsche Anreize zerstört werden (siehe Crowding-Out). \\
\midrule
Kahneman, Daniel & Israelisch-amerikanischer Psychologe, Nobelpreis für Wirtschaft 2002. Begründer der Verhaltensökonomie zusammen mit Amos Tversky. Hauptwerk: ``Thinking, Fast and Slow'' (2011). Zentrale Beiträge: Prospect Theory, Dual-System-Theorie, Heuristiken und Biases (Kahneman \& Tversky, 1979; Kahneman, 2011). \\
\midrule
Komplementarität ($\gamma$) & Wechselwirkung zwischen Faktoren, die in der Realität stattfinden (psychologisch, finanziell, kulturell). Verstärkend: Transparenz + Feedbackkultur wirken zusammen stärker. Neutralisierend: Individualboni + Teamziele heben sich auf. Blockierend: Kontrolle verhindert Vertrauenskultur. Im BCM als $\gamma$-Matrix quantifiziert (Kapitel 5). \\
\midrule
Kontext ($\Psi$) & Die acht Dimensionen, die bestimmen, ob eine Intervention wirkt: zeitlich (Dringlichkeit), sozial (Beobachtung), institutionell (Regeln), informationell (Wissen), emotional (Stress), ressourcenbezogen (Zeit/Geld), kulturell (Schweiz $\neq$ Schweden), physisch (Ort). Im BCM als $\Psi$-Vektor formalisiert (Kapitel 9). \\
\midrule
Libertärer Paternalismus & Philosophie hinter Nudging: Menschen zu besseren Entscheidungen führen, ohne ihre Wahlfreiheit einzuschränken. ``Libertär'' = Optionen bleiben offen. ``Paternalistisch'' = es gibt eine empfohlene Richtung (Thaler \& Sunstein, 2008). \\
\midrule
Loss Aversion & Menschen empfinden Verluste etwa doppelt so stark wie gleichgrosse Gewinne. CHF 100 zu verlieren schmerzt mehr, als CHF 100 zu gewinnen freut. Erklärt, warum Menschen irrationale Risiken eingehen, um Verluste zu vermeiden (Kahneman \& Tversky, 1979). \\
\midrule
Nudge & Verhaltensökonomische Intervention, die Entscheidungen beeinflusst, ohne Optionen zu verbieten oder finanzielle Anreize zu setzen. Nutzt psychologische Mechanismen wie Defaults, Framing oder soziale Normen. Begriff geprägt von Thaler \& Sunstein (2008). \\
\midrule
Present Bias & Tendenz, sofortige Belohnungen gegenüber zukünftigen zu bevorzugen -- selbst wenn die zukünftigen grösser wären. Erklärt Aufschieben, ungesunde Ernährung, mangelndes Sparen. Nicht Irrationalität, sondern menschlich. \\
\midrule
Prospect Theory & Theorie von Kahneman \& Tversky (1979), die erklärt, wie Menschen tatsächlich entscheiden. Kernaussagen: Verluste wiegen schwerer als Gewinne, Referenzpunkte bestimmen Bewertung, Menschen sind risikoscheu bei Gewinnen und risikofreudig bei Verlusten. \\
\midrule
Rationalität (begrenzte) & Konzept von Herbert Simon: Menschen sind nicht irrational, aber ihre Rationalität ist begrenzt durch Zeit, Information und kognitive Kapazität. Sie suchen nicht die optimale Lösung, sondern eine ``gut genuge'' (Satisficing). \\
\midrule
SAM & Serviceable Addressable Market. Der Teil des Gesamtmarktes, der mit dem eigenen Angebot tatsächlich erreichbar ist. Für BEATRIX: Organisations- und Verhaltensberatung im DACH-Raum, CHF 5--10 Milliarden. \\
\midrule
Social Proof & Menschen orientieren sich am Verhalten anderer, besonders bei Unsicherheit. ``Wenn alle das tun, wird es richtig sein.'' Beispiel: ``80\% der Hotelgäste verwenden ihre Handtücher mehrfach'' erhöht Wiederverwendung um 26\%. \\
\midrule
SOM & Serviceable Obtainable Market. Der realistisch erreichbare Marktanteil innerhalb eines definierten Zeitraums. Für BEATRIX: CHF 50--100 Millionen in fünf Jahren. \\
\midrule
SSBM & Social Sciences \& Business Management. Fachgebiet bei Innosuisse für Management- und Sozialwissenschaften -- im Gegensatz zu ICT (Informatik). BEATRIX wird bei SSBM eingereicht, weil die Innovation in der Beratungsmethodik liegt, nicht in der Software. \\
\midrule
Status Quo Bias & Tendenz, beim aktuellen Zustand zu bleiben, selbst wenn Alternativen objektiv besser wären. Ursachen: Gewohnheit, Verlustangst, Aufwand der Veränderung. Erklärt, warum Menschen bei schlechten Tarifen, Versicherungen oder Gewohnheiten bleiben. \\
\midrule
Sunk Cost Fallacy & Irrtum, vergangene Investitionen bei Entscheidungen zu berücksichtigen, obwohl sie nicht mehr rückgängig gemacht werden können. ``Ich habe schon so viel investiert, jetzt muss ich weitermachen.'' Rational wäre: Nur zukünftige Kosten und Nutzen zählen. \\
\midrule
TAM & Total Addressable Market. Der theoretische Gesamtmarkt für ein Produkt oder eine Dienstleistung, ohne Einschränkungen. Für BEATRIX: Globaler Beratungsmarkt, USD 300+ Milliarden. \\
\midrule
Thaler, Richard & Amerikanischer Ökonom, Nobelpreis 2017, Professor an der University of Chicago. Mitbegründer der Verhaltensökonomie, Erfinder des Begriffs ``Nudge''. Forschung zu Mental Accounting, Selbstkontrolle, Fairness (Thaler \& Sunstein, 2008; Thaler, 2015). \\
\midrule
TRL & Technology Readiness Level. Skala von 1--9 zur Bewertung des Entwicklungsstands. TRL 1--3: Forschung. TRL 4--6: Prototyp. TRL 7--9: Marktreif. BEATRIX steht aktuell bei TRL 4--5, Ziel nach Innovation Cheque: TRL 6--7. \\
\midrule
Verhaltensökonomie & Wirtschaftswissenschaft, die psychologische Erkenntnisse in ökonomische Modelle integriert. Zeigt, dass Menschen systematisch von ``rationalem'' Verhalten abweichen. Nobelpreise: Kahneman (2002), Thaler (2017). \\
\midrule
Willingness (WAX) & Handlungsbereitschaft einer Person. Im BCM (Kapitel 12) mit 99 Axiomen formalisiert. Hängt ab von: Awareness (weiss ich, was ich tun soll?), Motivation (will ich es?), und wahrgenommenen Barrieren (kann ich es?). Erklärt, warum Wissen allein selten zu Verhaltensänderung führt. \\
\midrule
9C-CORE & Die neun Kernfragen des BCM (abgeleitet aus dem EBF-Framework): WHO (wer hat Nutzen?), WHAT (was ist Nutzen?), HOW (wie interagieren Faktoren?), WHEN (wann zählt Kontext?), WHERE (woher die Zahlen?), AWARE (wie bewusst?), READY (handlungsbereit?), STAGE (wo in der Journey?), HIERARCHY (wie stratifizieren Entscheidungen?). \\
\bottomrule
\end{longtable}

\newpage

%% =============================================================================
%% LITERATURVERZEICHNIS
%% =============================================================================

\section*{Literaturverzeichnis}
\addcontentsline{toc}{section}{Literaturverzeichnis}

\subsection*{Literatur}

Benartzi, S., \& Thaler, R. H. (2004). Save More Tomorrow: Using behavioral economics to increase employee saving. \textit{Journal of Political Economy}, 112(S1), S164--S187.

\vspace{0.3cm}
Charness, G., \& Rabin, M. (2002). Understanding social preferences with simple tests. \textit{Quarterly Journal of Economics}, 117(3), 817--869.

\vspace{0.3cm}
Deci, E. L., Koestner, R., \& Ryan, R. M. (1999). A meta-analytic review of experiments examining the effects of extrinsic rewards on intrinsic motivation. \textit{Psychological Bulletin}, 125(6), 627--668.

\vspace{0.3cm}
DellaVigna, S. (2009). Psychology and economics: Evidence from the field. \textit{Journal of Economic Literature}, 47(2), 315--372.

\vspace{0.3cm}
Falk, A., \& Fischbacher, U. (2006). A theory of reciprocity. \textit{Games and Economic Behavior}, 54(2), 293--315.

\vspace{0.3cm}
Fehr, E., \& Gächter, S. (2000). Cooperation and punishment in public goods experiments. \textit{American Economic Review}, 90(4), 980--994.

\vspace{0.3cm}
Fehr, E., \& Gächter, S. (2002). Altruistic punishment in humans. \textit{Nature}, 415(6868), 137--140.

\vspace{0.3cm}
Fehr, E., \& Gächter, S. (2002). Do incentive contracts undermine voluntary cooperation? \textit{Journal of Political Economy}, 110(4), 817--838.

\vspace{0.3cm}
Fehr, E., \& Schmidt, K. M. (1999). A theory of fairness, competition, and cooperation. \textit{Quarterly Journal of Economics}, 114(3), 817--868.

\vspace{0.3cm}
Frederick, S., Loewenstein, G., \& O'Donoghue, T. (2002). Time discounting and time preference: A critical review. \textit{Journal of Economic Literature}, 40(2), 351--401.

\vspace{0.3cm}
Frey, B. S., \& Oberholzer-Gee, F. (1997). The cost of price incentives: An empirical analysis of motivation crowding-out. \textit{American Economic Review}, 87(4), 746--755.

\vspace{0.3cm}
Gneezy, U., Meier, S., \& Rey-Biel, P. (2011). When and why incentives (don't) work to modify behavior. \textit{Journal of Economic Perspectives}, 25(4), 191--210.

\vspace{0.3cm}
Gneezy, U., \& Rustichini, A. (2000). Pay enough or don't pay at all. \textit{Quarterly Journal of Economics}, 115(3), 791--810.

\vspace{0.3cm}
Gneezy, U., \& Rustichini, A. (2000). A fine is a price. \textit{Journal of Legal Studies}, 29(1), 1--17.

\vspace{0.3cm}
Johnson, E. J., \& Goldstein, D. (2003). Do defaults save lives? \textit{Science}, 302(5649), 1338--1339.

\vspace{0.3cm}
Kahneman, D. (2011). \textit{Thinking, fast and slow}. New York: Farrar, Straus and Giroux.

\vspace{0.3cm}
Kahneman, D., Knetsch, J. L., \& Thaler, R. H. (1990). Experimental tests of the endowment effect and the Coase theorem. \textit{Journal of Political Economy}, 98(6), 1325--1348.

\vspace{0.3cm}
Kahneman, D., Knetsch, J. L., \& Thaler, R. H. (1991). Anomalies: The endowment effect, loss aversion, and status quo bias. \textit{Journal of Economic Perspectives}, 5(1), 193--206.

\vspace{0.3cm}
Kahneman, D., \& Tversky, A. (1979). Prospect theory: An analysis of decision under risk. \textit{Econometrica}, 47(2), 263--291.

\vspace{0.3cm}
Kotter, J. P. (1996). \textit{Leading change}. Boston: Harvard Business School Press.

\vspace{0.3cm}
Krakowski, S., Pashkevich, N., Luger, J., Haftor, D., \& Raisch, S. (2025). Human-Centered Artificial Intelligence: A Field Experiment. \textit{Management Science}, forthcoming.

\vspace{0.3cm}
Laibson, D. (1997). Golden eggs and hyperbolic discounting. \textit{Quarterly Journal of Economics}, 112(2), 443--478.

\vspace{0.3cm}
Loewenstein, G. (1996). Out of control: Visceral influences on behavior. \textit{Organizational Behavior and Human Decision Processes}, 65(3), 272--292.

\vspace{0.3cm}
Madrian, B. C., \& Shea, D. F. (2001). The power of suggestion: Inertia in 401(k) participation and savings behavior. \textit{Quarterly Journal of Economics}, 116(4), 1149--1187.

\vspace{0.3cm}
McKinsey \& Company (2015). Changing change management: A blueprint that takes hold. \textit{McKinsey Quarterly}.

\vspace{0.3cm}
O'Donoghue, T., \& Rabin, M. (1999). Doing it now or later. \textit{American Economic Review}, 89(1), 103--124.

\vspace{0.3cm}
Rabin, M. (1993). Incorporating fairness into game theory and economics. \textit{American Economic Review}, 83(5), 1281--1302.

\vspace{0.3cm}
Samuelson, W., \& Zeckhauser, R. (1988). Status quo bias in decision making. \textit{Journal of Risk and Uncertainty}, 1(1), 7--59.

\vspace{0.3cm}
Sunstein, C. R. (2014). \textit{Why nudge? The politics of libertarian paternalism}. New Haven: Yale University Press.

\vspace{0.3cm}
Thaler, R. H. (1985). Mental accounting and consumer choice. \textit{Marketing Science}, 4(3), 199--214.

\vspace{0.3cm}
Thaler, R. H. (1999). Mental accounting matters. \textit{Journal of Behavioral Decision Making}, 12(3), 183--206.

\vspace{0.3cm}
Thaler, R. H. (2015). \textit{Misbehaving: The making of behavioral economics}. New York: W.W. Norton \& Company.

\vspace{0.3cm}
Thaler, R. H., \& Sunstein, C. R. (2008). \textit{Nudge: Improving decisions about health, wealth, and happiness}. New Haven: Yale University Press.

\vspace{0.3cm}
Tversky, A., \& Kahneman, D. (1974). Judgment under uncertainty: Heuristics and biases. \textit{Science}, 185(4157), 1124--1131.

\vspace{0.3cm}
Tversky, A., \& Kahneman, D. (1981). The framing of decisions and the psychology of choice. \textit{Science}, 211(4481), 453--458.

\vspace{2cm}

\begin{center}
\rule{0.4\textwidth}{0.5pt}\\[0.5em]
{\small Dokument erstellt: 27. Januar 2026 | Aktualisiert: 12. Februar 2026}\\
{\small Version: 4.7 -- Scharfe Abgrenzung EBF / BCM / BEATRIX}\\
{\small Für: Innosuisse Innovation Cheque (SSBM)}
\end{center}

\end{document}
